% !TeX program = pdflatex
\documentclass[aspectratio=169,xcolor={dvipsnames,table}]{beamer}

\input{../../_en_preambulo_beamer}
% Comandos y macros estadísticas compartidas para las presentaciones Beamer en inglés (Metropolis)

% Conjuntos numéricos básicos
\newcommand{\R}{\mathbb{R}}
\newcommand{\N}{\mathbb{N}}
\newcommand{\Q}{\mathbb{Q}}
\newcommand{\Z}{\mathbb{Z}}
\newcommand{\set}[1]{\left\{ #1 \right\}}
\newcommand{\abs}[1]{\left|#1\right|}
\newcommand{\norm}[1]{\left\| #1 \right\|}

% Comandos estadísticos y probabilísticos del curso
\newcommand{\Var}{\operatorname{Var}}
\newcommand{\cov}{\operatorname{Cov}}
\newcommand{\corr}{\rho}
\newcommand{\comb}[2]{\begin{pmatrix} #1 \\ #2 \end{pmatrix}}
\newcommand{\s}{\sigma}
\newcommand{\particion}{\mathfrak{P}}
\newcommand{\card}[1]{\#\left( #1 \right)}
\newcommand{\seq}[3]{\set{#1}_{#2}^{#3}}
\newcommand{\E}{\mathbb{E}}
\newcommand{\Prob}{\mathbb{P}}

% Entornos pedagógicos y cajas en inglés académico natural (soportando tanto nombres en inglés como en español)
\newenvironment{discussion}{%
  \begin{alertblock}{Discussion Question}%
}{%
  \end{alertblock}%
}
\newenvironment{discusion}{%
  \begin{alertblock}{Discussion Question}%
}{%
  \end{alertblock}%
}

\newenvironment{reflection}{%
  \begin{alertblock}{Classroom Reflection}%
}{%
  \end{alertblock}%
}
\newenvironment{reflexion}{%
  \begin{alertblock}{Classroom Reflection}%
}{%
  \end{alertblock}%
}

\newenvironment{paradox}{%
  \begin{alertblock}{Exploratory Paradox}%
}{%
  \end{alertblock}%
}
\newenvironment{paradoja}{%
  \begin{alertblock}{Exploratory Paradox}%
}{%
  \end{alertblock}%
}

\newenvironment{motivation}{%
  \begin{exampleblock}{Motivation: Data Science in Practice}%
}{%
  \end{exampleblock}%
}
\newenvironment{motivacion}{%
  \begin{exampleblock}{Motivation: Data Science in Practice}%
}{%
  \end{exampleblock}%
}

\newenvironment{quickchallenge}{%
  \begin{block}{Quick Team Challenge}%
}{%
  \end{block}%
}
\newenvironment{retoclase}{%
  \begin{block}{Quick Team Challenge}%
}{%
  \end{block}%
}


\title{Discrete Cumulative Distribution Function}
\subtitle{Section 03.02 --- Cumulative Distribution Functions (Discrete CDF)}
\author[J. Castillo Colmenares]{Juliho Castillo Colmenares}
\institute[Tec de Monterrey]{Tecnológico de Monterrey}
\date{}

\begin{document}

% SLIDE 01: Cover Page
\begin{frame}[plain]
	\titlepage
\end{frame}

% SLIDE 02: Roadmap
\begin{frame}{Roadmap --- Chapter 03: Discrete Random Variables}
	\begin{block}{Curricular Structure of Chapter 03}
		\small
		\begin{enumerate}
			\itemsep=0.04cm
			\item \textcolor{gray}{Section 03.01: Discrete Probability Mass Functions (PMF and Support)}
			\item \textbf{\textcolor{TecRojo}{Section 03.02: Cumulative Distribution Functions (Discrete CDF)}} $\leftarrow$ \emph{Current focus}
			\item \textcolor{gray}{Section 03.03: Mathematical Expectation, Variance, and Moments}
			\item \textcolor{gray}{Section 03.04: Bernoulli and Binomial Distributions}
			\item \textcolor{gray}{Section 03.05: Geometric and Negative Binomial Distributions}
			\item \textcolor{gray}{Section 03.06: Hypergeometric Distribution}
		\end{enumerate}
	\end{block}
	\vspace{-0.15cm}
	\begin{alertblock}{Didactic Objective of Section 03.02}
		\footnotesize
		Master the mathematical analysis of the Cumulative Distribution Function $F_X(x) = P(X \le x)$ for discrete variables, geometrically characterizing its flat plateaus, computing jumps of mass via finite differences ($\Delta F$), and solving probabilities of arbitrary intervals.
	\end{alertblock}
\end{frame}

% SLIDE 03: Motivation and Intuition
\begin{frame}{Motivation --- Why Accumulate Probability?}
	\begin{columns}[T]
		\begin{column}{0.48\textwidth}
			\begin{block}{Limitations of Point Queries}
				\small
				\begin{itemize}
					\itemsep=0.06cm
					\item The PMF $f(x) = P(X=x)$ answers only questions of \emph{exact equality}.
					\item In engineering applications, risk management, and reliability analysis, critical queries involve \textbf{thresholds} or \textbf{intervals}:
					\begin{itemize}
						\item \scriptsize What is the probability that a server fails within $t \le 5$ ms?
						\item \scriptsize What is the probability of having at most $3$ dropped packets?
					\end{itemize}
				\end{itemize}
			\end{block}
		\end{column}
		\begin{column}{0.48\textwidth}
			\begin{block}{The Cumulative Operator}
				\small
				\begin{itemize}
					\itemsep=0.06cm
					\item Instead of repeatedly summing point masses across vast domains, we define a global function $F(x)$ that stores the \emph{total probabilistic history} up to point $x$.
					\item This function transforms complex summation tasks into simple evaluations and algebraic subtractions at boundary points.
				\end{itemize}
			\end{block}
		\end{column}
	\end{columns}
\end{frame}

% SLIDE 05: Formal Definition of the CDF
\begin{frame}{Formal Definition of the Cumulative Distribution Function}
	\begin{block}{Analytical Definition (Equation 2.6.2 of Master Book)}
		\scriptsize
		The \textbf{Cumulative Distribution Function} (CDF) of a discrete random variable $X$ is defined for any real number $x \in \mathbb{R}$ as:
		\vspace{-0.15cm}
		\begin{align*}
			F(x) = P(X \le x) = \sum_{u \le x} f(u)
		\end{align*}
		\vspace{-0.25cm}
		where the summation extends over all elements of the support $u \in \mathcal{S}_X$ satisfying $u \le x$.
	\end{block}
	\vspace{-0.4cm}
	\begin{columns}[T]
		\begin{column}{0.48\textwidth}
			\begin{exampleblock}{Ordered Finite Support}
				\scriptsize
				If $x_1 < x_2 < \dots < x_n$, then:
				\vspace{-0.15cm}
				\begin{align*}
					F(x) = \begin{cases}
						0 & x < x_1 \\
						f(x_1) & x_1 \le x < x_2 \\
						f(x_1) + f(x_2) & x_2 \le x < x_3 \\
						\vdots & \vdots \\
						1 & x \ge x_n
					\end{cases}
				\end{align*}
			\end{exampleblock}
		\end{column}
		\begin{column}{0.48\textwidth}
			\begin{alertblock}{Geometric Interpretation}
				\scriptsize
				\begin{itemize}
					\itemsep=0.02cm
					\item $F(x)$ advances via \textbf{horizontal plateaus} where probability mass is absent.
					\item At each support point $x_k \in \mathcal{S}_X$, the function undergoes a \textbf{vertical jump} of exact height $f(x_k)$.
				\end{itemize}
			\end{alertblock}
		\end{column}
	\end{columns}
\end{frame}

% SLIDE 06: Fundamental Properties I (Monotonicity & Continuity)
\begin{frame}{Fundamental Properties --- Monotonicity and Right-Continuity}
	\begin{block}{Property 1: Non-Decreasing Monotonicity}
		\footnotesize
		For any two real numbers $a, b \in \mathbb{R}$ such that $a \le b$, it holds rigorously that:
		\begin{align*}
			a \le b \implies \set{\omega : X(\omega) \le a} \subseteq \set{\omega : X(\omega) \le b} \implies F(a) \le F(b)
		\end{align*}
		The accumulation of probabilities can never decrease, since point masses are non-negative ($f(x) \ge 0$).
	\end{block}
	\vspace{-0.15cm}
	\begin{block}{Property 2: Right-Continuity (cadlág)}
		\footnotesize
		At every point $x \in \mathbb{R}$, the right-hand limit coincides precisely with the evaluated function value:
		\begin{align*}
			F(x^+) = \lim_{u \to x^+} F(u) = F(x)
		\end{align*}
		At jump discontinuities $x_k \in \mathcal{S}_X$, the left-hand limit is strictly lower: $F(x_k^-) = F(x_k) - f(x_k)$. By universal mathematical convention, the upper point of the step is included (right-closure).
	\end{block}
\end{frame}

% SLIDE 07: Fundamental Properties II (Asymptotic Limits)
\begin{frame}{Fundamental Properties --- Asymptotic Limits and Normalization}
	\begin{block}{Property 3: Asymptotic Boundary Conditions}
		\small
		Every legitimate cumulative distribution function must satisfy two fundamental boundary conditions at infinity:
		\begin{align*}
			\lim_{x \to -\infty} F(x) = 0 \quad \text{and} \quad \lim_{x \to \infty} F(x) = 1
		\end{align*}
	\end{block}
	\vspace{-0.15cm}
	\begin{columns}[T]
		\begin{column}{0.48\textwidth}
			\begin{alertblock}{Lower Limit ($\to -\infty$)}
				\footnotesize
				As we move infinitely to the left on the real line, we leave behind all support points $\mathcal{S}_X$. Accumulated probability vanishes to zero (an impossible event).
			\end{alertblock}
		\end{column}
		\begin{column}{0.48\textwidth}
			\begin{exampleblock}{Upper Limit ($\to \infty$)}
				\footnotesize
				Advancing infinitely to the right encompasses the entire sample space $\Omega$. The total sum of all jumps converges exactly to unit normalization $1$.
			\end{exampleblock}
		\end{column}
	\end{columns}
\end{frame}

% SLIDE 08: Inversion Theorem (PMF Recovery)
\begin{frame}{Inversion Theorem --- Recovering the PMF from the CDF}
	\begin{block}{Finite Difference Theorem (Equation 2.6.4 of Master Book)}
		\small
		The probability mass function $f(x)$ can be uniquely reconstructed from the cumulative distribution function $F(x)$ via the exact magnitude of the jump discontinuity at point $x$:
		\begin{align*}
			f(x) = F(x) - \lim_{u \to x^-} F(u) \equiv F(x) - F(x^-)
		\end{align*}
	\end{block}
	\vspace{-0.15cm}
	\begin{columns}[T]
		\begin{column}{0.48\textwidth}
			\begin{exampleblock}{Ordered Discrete Support}
				\scriptsize
				If $\mathcal{S}_X = \set{x_1, x_2, \dots, x_n}$ with $x_1 < x_2 < \dots$:
				\begin{align*}
					f(x_1) &= F(x_1) - 0 = F(x_1), \\
					f(x_k) &= F(x_k) - F(x_{k-1}) \quad (k > 1).
				\end{align*}
			\end{exampleblock}
		\end{column}
		\begin{column}{0.48\textwidth}
			\begin{alertblock}{Points Outside the Support}
				\scriptsize
				If $x \notin \mathcal{S}_X$, the function $F(u)$ is locally constant around $x$ (lying strictly within a horizontal plateau). Therefore:
				\begin{align*}
					F(x) = F(x^-) \implies f(x) = 0.
				\end{align*}
			\end{alertblock}
		\end{column}
	\end{columns}
\end{frame}

% SLIDE 09: Probability Algebra over Intervals
\begin{frame}{Probability Algebra over Intervals}
	\begin{block}{Left-Open, Right-Closed Intervals}
		\footnotesize
		For any pair of real numbers $a, b \in \mathbb{R}$ with $a \le b$, the probability that the random variable falls within $(a, b]$ equals the exact difference of their accumulations:
		\begin{align*}
			P(a < X \le b) = P(X \le b) - P(X \le a) = F(b) - F(a)
		\end{align*}
	\end{block}
	\vspace{-0.15cm}
	\begin{block}{Complete Typology of Discrete Intervals}
		\scriptsize
		In discrete variables, strict vs. non-strict boundary inequalities alter the probability depending on whether mass exists at $a$ and $b$:
		\begin{table}
			\centering
			\begin{tabular}{l|l|l}
				\hline
				\textbf{Interval Type} & \textbf{Probability Expression} & \textbf{CDF Decomposition} \\
				\hline \hline
				Half-open $(a, b]$ & $P(a < X \le b)$ & $F(b) - F(a)$ \\
				\hline
				Closed $[a, b]$ & $P(a \le X \le b) = P(X=a) + P(a < X \le b)$ & $F(b) - F(a^-)$ \\
				\hline
				Open $(a, b)$ & $P(a < X < b) = P(a < X \le b) - P(X=b)$ & $F(b^-) - F(a)$ \\
				\hline
				Half-open $[a, b)$ & $P(a \le X < b) = P(a \le X \le b) - P(X=b)$ & $F(b^-) - F(a^-)$ \\
				\hline
			\end{tabular}
		\end{table}
	\end{block}
\end{frame}

% SLIDE 11: Python Lab --- Part 1: CDF Construction
\begin{frame}[fragile]{Computational Lab --- Part 1: Step CDF Construction}
	\begin{block}{Programmatic Probability Accumulation (\texttt{compute\_step\_cdf})}
		\scriptsize
		The following Python block sorts the vector support and applies cumulative summation (\texttt{np.cumsum}) to compute the exact plateau heights of $F(x)$.
	\end{block}
	\vspace{-0.1cm}
	\lstinputlisting[language=Python, firstline=9, lastline=21, basicstyle=\fontsize{6.3pt}{7.5pt}\ttfamily]{../../code/03_variables_aleatorias_discretas/03.02_discrete_cdf.py}
\end{frame}

% SLIDE 12: Python Lab --- Part 2: PMF Jump Inversion
\begin{frame}[fragile]{Computational Lab --- Part 2: Jump Inversion and PMF Recovery}
	\begin{block}{Finite Difference Algorithm (\texttt{extract\_pmf\_from\_cdf})}
		\scriptsize
		Using the backward difference operator (\texttt{np.diff} with \texttt{prepend=0.0}), the script extracts exact jump masses and verifies numerical normalization exactness.
	\end{block}
	\vspace{-0.1cm}
	\lstinputlisting[language=Python, firstline=23, lastline=35, basicstyle=\fontsize{6.3pt}{7.5pt}\ttfamily]{../../code/03_variables_aleatorias_discretas/03.02_discrete_cdf.py}
\end{frame}

% SLIDE 13: Python Lab --- Part 3: Interval Queries & Quantiles
\begin{frame}[fragile]{Computational Lab --- Part 3: Interval Queries and Quantiles}
	\begin{block}{Subset Evaluation and Percentile Lookup (\texttt{evaluate\_interval\_and\_quantiles})}
		\scriptsize
		The code leverages binary search (\texttt{np.searchsorted}) to evaluate interval probabilities via CDF subtractions and determines the median and 85th percentile using logical thresholds.
	\end{block}
	\vspace{-0.1cm}
	\lstinputlisting[language=Python, firstline=37, lastline=52, basicstyle=\fontsize{6.3pt}{7.5pt}\ttfamily]{../../code/03_variables_aleatorias_discretas/03.02_discrete_cdf.py}
\end{frame}

% SLIDE 14: Python Lab --- Part 4: Monte Carlo ECDF Simulation
\begin{frame}[fragile]{Computational Lab --- Part 4: Monte Carlo Simulation and ECDF}
	\begin{block}{Empirical Distribution Convergence (\texttt{simulate\_empirical\_ecdf})}
		\scriptsize
		We simulate $N=100,000$ double-dice rolls to construct the ECDF $F_n(x)$ and confirm via the Glivenko-Cantelli Theorem a supremum absolute deviation below $0.005$.
	\end{block}
	\vspace{-0.2cm}
	\lstinputlisting[language=Python, firstline=54, lastline=73, basicstyle=\fontsize{5.4pt}{6.5pt}\ttfamily]{../../code/03_variables_aleatorias_discretas/03.02_discrete_cdf.py}
\end{frame}

% SLIDE 15A: Classroom Exercise --- Tier 1: Fundamental (Statement)
\begin{frame}{Classroom Exercise --- Tier 1: Fundamental (1/2)}
	\begin{block}{Problem 3.2.1 --- Fair Die and Evaluation Points in $\mathbb{R}$}
		\small
		Consider a fair six-sided die whose outcome is the discrete random variable $X \in \{1, 2, 3, 4, 5, 6\}$ with uniform probabilities $P(X=k) = \frac{1}{6}$.
		\begin{enumerate}\itemsep=0.15cm
			\item Construct the exhaustive piecewise analytical expression of the cumulative distribution function $F_X(x) = P(X \le x)$ for all $x \in \mathbb{R}$.
			\item Compute and interpret the exact numerical value of the function at the following real points of the domain: $x = 2.8$, $x = -1.0$, $x = 4.0$, and $x = 10.0$.
		\end{enumerate}
	\end{block}
\end{frame}

% SLIDE 15B: Classroom Exercise --- Tier 1: Fundamental (Solution)
\begin{frame}{Classroom Exercise --- Tier 1: Fundamental (2/2)}
	\fontsize{6.8pt}{8.2pt}\selectfont
	\begin{block}{1. Construction of the Piecewise Analytical Expression ($F_X(x)$)}
		\fontsize{6.8pt}{8.2pt}\selectfont
		Progressively accumulating jumps of magnitude $\frac{1}{6}$ at each integer of the support:
		$$F_X(x) = \begin{cases}
			0 & \text{if } x < 1 \\
			1/6 \approx 0.1667 & \text{if } 1 \le x < 2 \\
			2/6 \approx 0.3333 & \text{if } 2 \le x < 3 \\
			3/6 = 0.5000 & \text{if } 3 \le x < 4 \\
			4/6 \approx 0.6667 & \text{if } 4 \le x < 5 \\
			5/6 \approx 0.8333 & \text{if } 5 \le x < 6 \\
			1 & \text{if } x \ge 6
		\end{cases}$$
	\end{block}
	\vspace{-0.15cm}
	\pause
	\begin{alertblock}{2. Point Evaluation at Real Arguments $\mathbb{R}$}
		\fontsize{6.8pt}{8.2pt}\selectfont
		\begin{itemize}\itemsep=0.06cm
			\item $F_X(2.8):$ Since $2 \le 2.8 < 3$, it falls in the second plateau $\implies \mathbf{F_X(2.8) = \frac{2}{6} = \frac{1}{3} \approx 0.3333}$.
			\item $F_X(-1.0):$ Since $-1.0 < 1$ (lower boundary of the support), it lies to the left $\implies \mathbf{F_X(-1.0) = 0}$.
			\item $F_X(4.0):$ Since the interval includes the left endpoint ($4 \le 4.0 < 5$) $\implies \mathbf{F_X(4.0) = \frac{4}{6} = \frac{2}{3} \approx 0.6667}$.
			\item $F_X(10.0):$ Since $10.0 \ge 6$ exceeds the entire support $\implies \mathbf{F_X(10.0) = 1.0000}$.
		\end{itemize}
	\end{alertblock}
\end{frame}

% SLIDE 16A: Classroom Exercise --- Tier 2: Operational (Statement)
\begin{frame}{Classroom Exercise --- Tier 2: Operational (1/2)}
	\begin{block}{Problem 3.2.4 --- Solder Defects in Circuit Wafers}
		\footnotesize
		On an automated integrated-circuit production line, the number of defective solder joints per wafer, $X$, has the following probability mass function:
		$$p_X(x) = \begin{cases}
			0.40 & \text{if } x = 0 \\
			0.30 & \text{if } x = 1 \\
			0.20 & \text{if } x = 2 \\
			0.10 & \text{if } x = 3 \\
			0 & \text{otherwise}
		\end{cases}$$
		\begin{enumerate}\itemsep=0.04cm
			\item Write the cumulative distribution function $F_X(x)$ explicitly, piece by piece, for all $x \in \mathbb{R}$.
			\item The quality control department's protocol establishes that a wafer is \textbf{formally rejected if it presents $2$ or more defects}. Compute the rejection probability operating \emph{exclusively} with values of the cumulative function $F_X(x)$.
		\end{enumerate}
	\end{block}
\end{frame}

% SLIDE 16B: Classroom Exercise --- Tier 2: Operational (Solution)
\begin{frame}{Classroom Exercise --- Tier 2: Operational (2/2)}
	\begin{columns}[T]
		\begin{column}{0.48\textwidth}
			\begin{block}{1. Piecewise Expression of $F_X(x)$}
				\fontsize{6.5pt}{7.8pt}\selectfont
				Summing point probabilities ($0.40 + 0.30 = 0.70$, etc.):
				$$F_X(x) = \begin{cases}
					0.00 & \text{if } x < 0 \\
					0.40 & \text{if } 0 \le x < 1 \\
					0.70 & \text{if } 1 \le x < 2 \\
					0.90 & \text{if } 2 \le x < 3 \\
					1.00 & \text{if } x \ge 3
				\end{cases}$$
			\end{block}
		\end{column}
		\begin{column}{0.48\textwidth}
			\begin{alertblock}{2. Rejection Probability ($P(X \ge 2)$)}
				\fontsize{6.5pt}{7.8pt}\selectfont
				Upper tail $X \ge 2$. By the complement axiom:
				$$P(\text{Rejection}) = P(X \ge 2) = 1 - P(X < 2)$$
				Since $X \in \mathbb{Z}$, the event $\{X < 2\}$ is equivalent to $\{X \le 1\}$:
				$$P(X \ge 2) = 1 - P(X \le 1) = 1 - F_X(1)$$
				Substituting the corresponding plateau of $F_X(x)$:
				$$P(\text{Rejection}) = 1 - 0.70 = \mathbf{0.30 \implies 30.0\%}$$
			\end{alertblock}
		\end{column}
	\end{columns}
\end{frame}

% SLIDE 17A: Classroom Exercise --- Tier 3: Analytical (Statement)
\begin{frame}{Classroom Exercise --- Tier 3: Analytical (1/2)}
	\begin{block}{Problem 3.2.7 --- Probability Theorem over Integer Intervals}
		\footnotesize
		Let $X$ be a discrete random variable whose support is a subset of the integers ($\mathcal{S}_X \subset \mathbb{Z}$) and whose cumulative distribution function is $F_X(x)$.
		\begin{enumerate}\itemsep=0.15cm
			\item Formally prove that for any two integers $a, b \in \mathbb{Z}$ with $a \le b$, the probability of the closed interval is expressed in terms of the CDF by:
			$$P(a \le X \le b) = F_X(b) - F_X(a - 1)$$
			\item From the previous result, rigorously prove that the magnitude of the jump at any integer $k \in \mathbb{Z}$ equals the point mass:
			$$p_X(k) = F_X(k) - F_X(k-1)$$
		\end{enumerate}
	\end{block}
\end{frame}

% SLIDE 17B: Classroom Exercise --- Tier 3: Analytical (Solution)
\begin{frame}{Classroom Exercise --- Tier 3: Analytical (2/2)}
	\fontsize{6.5pt}{7.8pt}\selectfont
	\begin{block}{1. Proof for the Closed Interval over $\mathbb{Z}$}
		\fontsize{6.5pt}{7.8pt}\selectfont
		By the axiomatic definition of the CDF, the accumulated event $\{X \le b\}$ can be partitioned into the union of two disjoint events: $\{X \le a-1\} \cup \{a \le X \le b\}$. By additivity:
		$$P(X \le b) = P(X \le a-1) + P(a \le X \le b) \implies P(a \le X \le b) = P(X \le b) - P(X \le a-1)$$
		Substituting the formal functional notation of the cumulative distribution ($F_X(x) = P(X \le x)$):
		$$P(a \le X \le b) = F_X(b) - F_X(a-1) \quad \blacksquare$$
	\end{block}
	\vspace{-0.25cm}
	\pause
	\begin{alertblock}{2. Proof of Point-Jump Inversion over $\mathbb{Z}$}
		\fontsize{6.5pt}{7.8pt}\selectfont
		Let $a = b = k$ be an arbitrary integer. Substituting both endpoints by $k$ in the previous closed-interval formula:
		$$P(k \le X \le k) = F_X(k) - F_X(k-1)$$
		The left-hand side represents the singleton event where the variable takes the exact value $k$, that is, $P(X=k) = p_X(k)$:
		$$p_X(k) = F_X(k) - F_X(k-1) = \Delta F_X(k) \quad \blacksquare$$
	\end{alertblock}
\end{frame}

% SLIDE 18A: Classroom Exercise --- Tier 4: Challenging (Statement)
\begin{frame}{Classroom Exercise --- Tier 4: Challenging (1/2)}
	\begin{block}{Problem 3.2.9 --- Spectral Representation via Heaviside Steps}
		\footnotesize
		In measure theory and probabilistic functional analysis, every pure discrete random variable can be represented as an infinite superposition of unit jumps.
		\begin{enumerate}\itemsep=0.12cm
			\item Prove that the cumulative distribution function $F_X(x)$ can be expressed as a convergent linear combination of Heaviside unit steps $H(x - x_k)$:
			$$F_X(x) = \sum_{k} p_X(x_k) H(x - x_k), \quad \text{where } H(t) = \begin{cases} 1 & \text{if } t \ge 0 \\ 0 & \text{if } t < 0 \end{cases}$$
			\item Explain why the ordinary derivative satisfies $F_X'(x) = 0$ almost everywhere (a.e.) with respect to Lebesgue measure, and how the theory of generalized distributions recovers the analytical density $f_X(x)$ via Dirac Delta impulses ($\delta(x-x_k)$).
		\end{enumerate}
	\end{block}
\end{frame}

% SLIDE 18B: Classroom Exercise --- Tier 4: Challenging (Solution)
\begin{frame}{Classroom Exercise --- Tier 4: Challenging (2/2)}
	\fontsize{6.5pt}{7.8pt}\selectfont
	\begin{block}{1. Proof of the Heaviside Series}
		\fontsize{6.5pt}{7.8pt}\selectfont
		We evaluate the series at an arbitrary point $x \in \mathbb{R}$. The step function $H(x-x_k)$ equals $1$ if and only if $x - x_k \ge 0 \iff x_k \le x$, and equals $0$ otherwise. Therefore, the infinite sum filters exactly the support masses to the left of $x$:
		$$\sum_{k} p_X(x_k) H(x - x_k) = \sum_{k: x_k \le x} p_X(x_k)(1) + \sum_{k: x_k > x} p_X(x_k)(0) = \sum_{x_k \le x} p_X(x_k) = F_X(x) \quad \blacksquare$$
	\end{block}
	\vspace{-0.15cm}
	\pause
	\begin{alertblock}{2. Generalized Derivative and the Dirac Delta}
		\fontsize{6.5pt}{7.8pt}\selectfont
		\begin{itemize}\itemsep=0.06cm
			\item At every point $x \notin \mathcal{S}_X$, $F_X(x)$ is locally constant $\implies F_X'(x) = 0$ almost everywhere (a.e.). The ordinary derivative loses all probability-mass information.
			\item In the space of tempered distributions ($\mathcal{D}'$), the derivative of the Heaviside step is the Dirac Delta: $\frac{d}{dx}H(x-x_k) = \delta(x-x_k)$. Differentiating term by term yields the impulse density:
			$$f_X(x) = \frac{d}{dx}F_X(x) = \sum_{k} p_X(x_k) \delta(x - x_k)$$
			whose integral over any interval $[a,b]$ exactly recovers the sum of discrete masses.
		\end{itemize}
	\end{alertblock}
\end{frame}

% SLIDE 19: Conclusions and Synthesis
\begin{frame}{Conclusions --- The CDF as a Pillar of Discrete Analysis}
	\begin{block}{Summary of Acquired Competencies}
		\small
		\begin{enumerate}
			\itemsep=0.06cm
			\item \textbf{Structural Equivalence:} The Cumulative Distribution Function $F(x)$ captures the entirety of an aleatory experiment with the exact same uniqueness and rigor as the PMF $f(x)$.
			\item \textbf{Step Geometry:} For discrete systems, $F(x)$ is characterized by flat plateaus interspersed with exact vertical jumps at support points $\mathcal{S}_X$.
			\item \textbf{Inversion and Algebra:} Backward finite differences $f(x) = F(x) - F(x^-)$ and interval formulas evaluate risk without tedious summation loops.
		\end{enumerate}
	\end{block}
	\vspace{-0.15cm}
	\begin{alertblock}{Computational Verification}
		\footnotesize
		Monte Carlo simulations confirmed that empirical distributions (ECDF) converge uniformly to theoretical step CDFs, validating the Fundamental Glivenko-Cantelli Theorem.
	\end{alertblock}
\end{frame}

% SLIDE 20: Didactic Bridge
\begin{frame}{Didactic Bridge --- Toward Mathematical Expectation and Moments}
	\begin{columns}[T]
		\begin{column}{0.48\textwidth}
			\begin{block}{What We Have Achieved (03.01 \& 03.02)}
				\small
				\begin{itemize}
					\itemsep=0.06cm
					\item We mastered the complete micro-analytical description of discrete variables via $f(x)$ (point masses) and $F(x)$ (step plateaus).
					\item We can answer exact, interval, quantile, and tail queries with precision.
				\end{itemize}
			\end{block}
		\end{column}
		\begin{column}{0.48\textwidth}
			\begin{block}{What Lies Ahead: Section 03.03}
				\small
				\begin{itemize}
					\itemsep=0.06cm
					\item How can we summarize the entire distribution into a single representative number of central tendency and dispersion?
					\item We will introduce the \textbf{Expected Value} operator $\mathbb{E}[X] = \sum x f(x)$, the \textbf{Variance} $\text{Var}(X)$, and the theory of moments and generating functions.
				\end{itemize}
			\end{block}
		\end{column}
	\end{columns}
	\vspace{0.2cm}
	\begin{center}
		\textbf{\textcolor{TecAzul}{\Large Congratulations on completing Section 03.02!}}
	\end{center}
\end{frame}

\end{document}
